\documentclass[12pt]{article}
\usepackage[utf8]{inputenc}
\usepackage[T2A]{fontenc}
\usepackage[russian]{babel}
\usepackage{amsmath,amssymb,amsthm}
\usepackage[a4paper,left=25mm,right=20mm,top=22mm,bottom=25mm]{geometry}
\usepackage{array}

\newtheorem{theorem}{Теорема}
\newtheorem{lemma}{Лемма}
\newtheorem{proposition}{Предложение}
\newtheorem{corollary}{Следствие}
\theoremstyle{definition}
\newtheorem{definition}{Определение}
\theoremstyle{remark}
\newtheorem{remark}{Замечание}

\newcommand{\R}{\mathbb{R}}
\newcommand{\Z}{\mathbb{Z}}
\newcommand{\E}{\mathbb{E}}
\newcommand{\diam}{\operatorname{diam}}
\newcommand{\dist}{\operatorname{dist}}

\title{Верхняя оценка $\chi(\R^4)\le 48$ для хроматического числа
четырёхмерного евклидова пространства%
\thanks{Черновик от 22.07.2026; аффилиация и данные автора подлежат уточнению.}}
\author{Леонид Иванов}
\date{22 июля 2026 г.}

\begin{document}
\maketitle

\begin{abstract}
Построена решёточная раскраска пространства $\R^4$ в $48$ цветов, при которой
одноцветные точки не реализуют ни одного расстояния из замкнутого отрезка
$[1,\ell]$ при любом $\ell\le 1{,}0396$. В частности, $\chi(\R^4)\le 48$: это
улучшает верхнюю оценку $\chi(\R^4)\le 49$, известную с 2006~г.
\cite{Ivanov2006,CoulsonPayne2007,ABPR}, и даёт отрицательный ответ на вопрос~1
работы Армана, Бондаренко, Прымака и Радченко \cite{ABPR} в части размерности
$n=4$: минимум числа цветов решёточных (sublattice) раскрасок $\E^4$ равен не
$49$, а не превосходит $48$. Раскраска задаётся явной рациональной решёткой
общего положения (ячейка Вороного --- комбинаторный пермутоэдр с $30$ фасетами
и $f$-вектором $(120,240,150,30)$), найденной численной оптимизацией метрики
по пространству квадратичных форм.
Результат сертифицирован исчерпывающим перебором всех $472\,440$ подрешёток
индекса $48$ и согласием трёх независимых геометрических алгоритмов с точностью
$10^{-9}$. Дополнительно вычислены точные нормированные ширины запрещённых
интервалов всех конструкций работы \cite{ABPR} в размерностях $4\le n\le 9$.
\end{abstract}

\section{Введение}

Хроматическим числом $\chi(\R^n)$ называется наименьшее число цветов,
достаточное для такой раскраски точек $\R^n$, что любые две точки на расстоянии
$1$ окрашены различно (проблема Нельсона--Хадвигера; см. обзоры
\cite{Raigorodskii2000,Soifer}). Более общо, для отрезка $[1,\ell]$,
$\ell\ge 1$, через $\chi\bigl(\R^n,[1,\ell]\bigr)$ обозначим наименьшее число
цветов раскраски, при которой одноцветные точки не реализуют ни одного
расстояния из $[1,\ell]$ (хроматическое число с интервалом запрещённых
расстояний; систематическое изучение начато в \cite{Ivanov2011}).

Известные границы в малых размерностях:
$5\le\chi(\R^2)\le 7$ \cite{deGrey,Soifer},
$6\le\chi(\R^3)\le 15$ \cite{Nechushtan,Coulson},
$9\le\chi(\R^4)$ \cite{ExooIsmailescu}. Верхняя оценка в размерности $4$
имеет следующую историю: $\chi(\R^4)\le 54$ (Радоичич, Тот \cite{RadoicicToth},
линейная раскраска на решётке $A_4^*$); $\chi(\R^4)\le 49$ --- анонсировано в
заметке \cite{Ivanov2006} и упомянуто в \cite{CoulsonPayne2007}; полное
доказательство, вместе с исчерпывающей проверкой того, что $49$ ---
минимальный индекс допустимой подрешётки для решётки $D_4$ (а $54$ --- для
$A_4^*$), дано Арманом, Бондаренко, Прымаком и Радченко \cite{ABPR}. Там же
(Open question~1) авторы сформулировали ожидание, что улучшение невозможно:
\emph{<<We believe that the answer is negative for $n=4,5$, i.e.
$\min\chi_s(\E^4,\Lambda,\Lambda',\ell)=49$~\ldots, where the minima are taken
over all full rank lattices $\Lambda$, sublattices $\Lambda'$ and all
$\ell>0$>>}.

Основной результат настоящей работы опровергает это ожидание.

\begin{theorem}\label{thm:main}
Пусть $\ell_0=1{,}0396$. Тогда
\[
\chi\bigl(\R^4,[1,\ell]\bigr)\;\le\;48
\qquad\text{при всех }1\le\ell\le\ell_0 .
\]
В частности, $\chi(\R^4)\le 48$.
\end{theorem}

Раскраска в теореме~\ref{thm:main} --- решёточная (классы смежности подрешётки
индекса $48$ раскрашивают ячейки Вороного явно заданной рациональной решётки
$\Lambda_\star\subset\R^4$), т.е. в терминологии \cite{ABPR}
$\min_{\Lambda,\Lambda',\ell}\chi_s(\E^4,\Lambda,\Lambda',\ell)\le 48$.
Существенно, что решётка $\Lambda_\star$ --- \emph{общего положения}
(генерическая): у неё лишь одна пара минимальных векторов, а ячейка Вороного
имеет максимально возможные $30$ фасет. Все ранее рассматривавшиеся
конструкции в $\R^4$ опирались на классические решётки ($D_4$, $A_4^*$,
$\Z^4$); исчерпывающие переборы работы \cite{ABPR} покрывали только $D_4$ и
$A_4^*$ и потому не могли обнаружить $\Lambda_\star$.

Отметим, что нижняя оценка Кулсона для решёточных раскрасок,
$\chi_s(\E^n,\Lambda,\Lambda',\ell)\ge 2^{n+1}-1$ \cite{Coulson,ABPR}, в
размерности $4$ равна $31$ и не препятствует значению $48$. Проверка значений
$k=47$ и $k=46$ той же оптимизационной схемой (раздел~\ref{sec:comp}) дала
максимумы нормированной ширины $0{,}973$ и $0{,}981<1$ соответственно, что
позволяет предположить, что для решёточных раскрасок $\R^4$ значение $48$
окончательно.

\section{Метод решёточных раскрасок}\label{sec:method}

Напомним конструкцию (Л.\,Л.~Иванов \cite{Ivanov2011}; в эквивалентной форме
--- \cite[предложение~5]{ABPR}). Пусть $\Lambda\subset\R^n$ --- решётка полного
ранга с ячейкой Вороного
$V_0=\{x:\ \|x\|\le\|x-v\|\ \forall v\in\Lambda\}$, и пусть
$\Gamma\subseteq\Lambda$ --- подрешётка индекса $k=[\Lambda:\Gamma]$. Раскрасим
ячейку $V_\Lambda(v)$, $v\in\Lambda$, цветом класса смежности $v+\Gamma$ ---
получится $k$-цветная периодическая раскраска $\R^n$ (границы ячеек
распределяются полуоткрытым образом, см. замечание~\ref{rem:boundary}).

\begin{lemma}[Вороной; см. \cite{Voronoi}]\label{lem:sym}
Ячейка $V_0$ --- выпуклый центрально-симметричный многогранник, и все её грани
центрально-симметричны.
\end{lemma}

\begin{lemma}[лемма о середине; {\cite[лемма~2]{Ivanov2011}}]\label{lem:mid}
Минимальное расстояние между ячейками $V_\Lambda(a)$ и $V_\Lambda(b)$
реализуется на паре точек, симметричных относительно середины отрезка $[a,b]$;
как следствие, для $v\in\Lambda$
\[
D(v):=\dist\bigl(V_0,\,v+V_0\bigr)=2\,\dist\bigl(v/2,\,V_0\bigr).
\]
\end{lemma}

Положим
\[
D(\Gamma)=\min_{v\in\Gamma\setminus\{0\}}D(v),\qquad
d(\Lambda,\Gamma)=\frac{D(\Gamma)}{\diam V_0}.
\]
Расстояния между одноцветными точками лежат в
$[0,\diam V_0)\cup[D(\Gamma),+\infty)$, поэтому после подходящей гомотетии
раскраска запрещает весь отрезок $[1,\ell]$ при любом
$\ell\le d(\Lambda,\Gamma)$:

\begin{proposition}[{ср. \cite[Prop.~5]{ABPR}}]\label{prop:crit}
Если $d(\Lambda,\Gamma)\ge 1$, то
$\chi\bigl(\R^n,[1,\ell]\bigr)\le[\Lambda:\Gamma]$ при всех
$\ell\le d(\Lambda,\Gamma)$.
\end{proposition}

\begin{remark}\label{rem:boundary}
При $\ell<d(\Lambda,\Gamma)$ строгих неравенств хватает с запасом: выберем
масштаб $s$ так, что $s\,\diam V_0<1$ и $s\,D(\Gamma)>\ell$; тогда никакие
граничные тонкости не возникают. Случай $\ell=d$ обслуживается стандартным
полуоткрытым распределением границ (как у Кулсона \cite{Coulson} и в
\cite[доказательство Prop.~5]{ABPR}). Для теоремы~\ref{thm:main} достаточно
первого случая, так как $\ell_0$ выбрано строго меньше $d$.
\end{remark}

Практическое вычисление $D(\Gamma)$ опирается на неравенство
$D(v)\ge\|v\|-\diam V_0$ (ячейка лежит в шаре радиуса $\tfrac12\diam V_0$):
достаточно перебрать конечное множество векторов
$\{v\in\Gamma:\ \|v\|<D_{\text{тек}}+\diam V_0\}$, где $D_{\text{тек}}$ ---
текущая верхняя оценка минимума; перебор полон при любом выборе базиса.

\section{Конструкция}\label{sec:construction}

Определим решётку $\Lambda_\star=B^{\mathsf T}\Z^4$, где $B$ --- верхний
треугольный фактор Холецкого ($Q=BB^{\mathsf T}$, строки $B$ --- базисные
векторы) следующей матрицы Грама с знаменателями, не превосходящими $12$:
\begin{equation}\label{eq:gram}
Q=\begin{pmatrix}
\tfrac{137}{12} & \tfrac{13}{7}  & -\tfrac{61}{10} & -\tfrac{18}{5}\\[2pt]
\tfrac{13}{7}   & 7              & -\tfrac{25}{9}  & -\tfrac{37}{12}\\[2pt]
-\tfrac{61}{10} & -\tfrac{25}{9} & 12              & -\tfrac{17}{11}\\[2pt]
-\tfrac{18}{5}  & -\tfrac{37}{12}& -\tfrac{17}{11} & \tfrac{44}{5}
\end{pmatrix},
\qquad
13860\,Q\in\mathrm{Mat}_4(\Z).
\end{equation}
Подрешётку $\Gamma_\star\subset\Lambda_\star$ индекса $48$ зададим матрицей
перехода в эрмитовой нормальной форме (строки коэффициентов относительно
базиса $B$):
\begin{equation}\label{eq:hnf}
T=\begin{pmatrix}
1&0&0&28\\
0&1&0&16\\
0&0&1&30\\
0&0&0&48
\end{pmatrix},
\qquad \det T=48 .
\end{equation}

Числовые инварианты (все значения устойчиво воспроизводятся тремя независимыми
алгоритмами, см. раздел~\ref{sec:comp}; величины $\lambda_1^2$, $D^2$,
$\diam^2 V_0$ рациональны, поскольку рациональны $Q$ и, следовательно, все
вершины ячейки):
\[
\lambda_1(\Lambda_\star)^2=7,\qquad
\diam^2 V_0\approx 19{,}1416,\qquad
D(\Gamma_\star)^2\approx 20{,}6899,
\]
\begin{equation}\label{eq:d}
d(\Lambda_\star,\Gamma_\star)=\frac{D(\Gamma_\star)}{\diam V_0}
=1{,}039657681\ldots>\ell_0=1{,}0396 .
\end{equation}
Ячейка Вороного $V_0(\Lambda_\star)$ имеет $30$ фасет (15 пар релевантных
векторов --- максимум в размерности~4) и $f$-вектор $(120,240,150,30)$
комбинаторного пермутоэдра; у $\Lambda_\star$ единственная пара минимальных
векторов, отношение покрытие/упаковка $2R/\lambda_1\approx1{,}577$. Решётка не
изометрична ни одной из классических ($D_4$, $A_4^*$, $\Z^4$, $A_4$).

Теорема~\ref{thm:main} немедленно следует из
предложения~\ref{prop:crit} и неравенства \eqref{eq:d}.

\section{Вычислительное доказательство неравенства \eqref{eq:d}}\label{sec:comp}

Величина $d(\Lambda_\star,\Gamma_\star)$ вычислена тремя независимыми
программами; все три дают $1{,}039657681$ с согласием в девяти значащих
цифрах.

\textbf{(i) Сертифицированное построение ячейки.} Релевантные векторы находятся
точным CVP-перебором в каждом из $2^4-1$ нетривиальных классов смежности
$\Lambda/2\Lambda$ (сдвинутое сферическое декодирование со стартовым радиусом
Бабаи; вектор релевантен тогда и только тогда, когда минимум класса единственен
с точностью до знака --- теорема Вороного). Построенная ячейка проходит
контроль: соотношение Эйлера для $f$-вектора, центральная симметрия множества
вершин, равенство $\operatorname{vol}V_0=\det\Lambda_\star$ (проверено
независимой библиотекой выпуклых оболочек с точностью $10^{-9}$).

\textbf{(ii) Три способа вычисления $\dist(v/2,V_0)$.} (а) алгоритм
Гилберта--Джонсона--Кирти по вершинам ячейки с апостериорным сертификатом
оптимальности через зазор двойственности (не более $10^{-6}$ от значения);
(б) каскад ортогональных проекций по граням (независимая реализация);
(в) решение задачи квадратичного программирования
$\min\|x-v/2\|^2$ при $Ax\le b$ по полупространствам релевантных векторов.

\textbf{(iii) Полнота перебора.} Минимум $D(\Gamma_\star)$ берётся по всем
векторам $\Gamma_\star$ в доказуемо достаточном шаре
$\|v\|<D_{\text{тек}}+\diam V_0$ (раздел~\ref{sec:method}). Кроме того,
выполнен \emph{исчерпывающий} перебор всех подрешёток индекса $48$ решётки
$\Lambda_\star$: их ровно
$\sum_{d_1d_2d_3d_4=48}d_2d_3^2d_4^3=472\,440$ (перечисление по эрмитовым
нормальным формам), и максимум $d$ достигается на $\Gamma_\star$ ---
т.е. \eqref{eq:hnf} оптимальна для $\Lambda_\star$ на индексе $48$.

\textbf{(iv) Направление погрешностей.} Возможные источники неточности
построения смещают вычисленное $d$ \emph{вниз}, а не вверх: пропуск
релевантного вектора раздувает ячейку (уменьшая $D$ и увеличивая знаменатель),
а алгоритм (а) возвращает верхнюю оценку расстояния, контролируемую
сертификатом. Запас $d-1>0{,}039$ на четыре порядка превосходит суммарные
допуски ($\le 10^{-6}$).

\textbf{(v) Как найдена решётка.} Максимизировалась функция
$t\mapsto d\bigl(Q(t),48\bigr)$ методом Нелдера--Мида: сначала по
$10$-мерному вторичному конусу генерического типа Делоне (в обозначениях
\cite{VWZ} --- тип $111{-}$; у его стандартного представителя
$d(48)=0{,}99941$), затем по всем положительно определённым формам в
параметризации Холецкого. Численный оптимум $d(48)=1{,}04327$ достигается на
иррациональной форме; матрица \eqref{eq:gram} --- ближайшая форма с малыми
знаменателями, сохраняющая $d>\ell_0$. Скрипты поиска и верификации, а также
все промежуточные данные доступны у автора (пакеты \texttt{voronoi4d} и
\texttt{combigeo}, версия~1.1.0).

\section{Сопутствующие результаты}\label{sec:extra}

Тем же инструментарием вычислены \emph{точные} нормированные ширины
запрещённых интервалов всех конструкций работы \cite{ABPR} (в самой работе
доказывалось только $d\ge1$); каждая строка даёт оценку
$\chi(\R^n,[1,\ell])\le k$ при $\ell\le d$:

\begin{center}
\begin{tabular}{c c l l}
\hline
$n$ & $k$ & конструкция & $d$ (точно)\\
\hline
4 & 49  & $(3+\omega)D_4$ & $\sqrt{7/6}=1{,}080123\ldots$\\
4 & 54  & $A_4^*$ & $\sqrt{23/20}=1{,}072381\ldots$\\
5 & 140 & $A_5^*$, матрица $C_5$ из \cite{ABPR} & $\sqrt{39/35}=1{,}055597\ldots$\\
6 & 343 & $(3+\omega)E_6^*$ & $\sqrt{7/6}$\\
7 & 1372& $E_7^*$, матрица $C_7$ из \cite{ABPR} & $\sqrt{17/14}=1{,}101946\ldots$\\
8 & 2401& $(3+\omega)E_8$ & $\sqrt{7/6}$\\
9 & 17253& $A_9^*$, матрица $C_9$ из \cite{ABPR} & $13/\sqrt{165}=1{,}012048\ldots$\\
\hline
\end{tabular}
\end{center}

Для эйзенштейновской конструкции $\Gamma=(3+\omega)\Lambda$ (умножение на
$\alpha=3+\omega$, $|\alpha|^2=7$, в решётках с комплексной структурой над
$\Z[\omega]$) во всех проверенных случаях ($n=2,4,6,8$) выполняется точное
равенство
\[
D\bigl((3+\omega)\Lambda\bigr)^2=\tfrac{7}{3}\,\lambda_1(\Lambda)^2,
\qquad\text{т.е.}\qquad
d=\frac{\sqrt{7/3}}{2R/\lambda_1},
\]
где $R$ --- радиус покрытия; в частности, при $2R/\lambda_1=\sqrt2$ (решётки
$D_4$, $E_6^*$, $E_8$, $\Lambda_{24}$) всегда $d=\sqrt{7/6}$. Отсюда, в
предположении сохранения равенства при $n=24$, следует гипотетическая оценка
$\chi\bigl(\R^{24},[1,\sqrt{7/6}]\bigr)\le 7^{12}$. Для решётки
Кокстера--Тодда $K_{12}$ ($2R/\lambda_1=\sqrt{8/3}$) та же формула даёт
$d\approx0{,}935<1$: конструкция теоремы~3 работы \cite{ABPR} на $K_{12}$ не
проходит (частичный отрицательный ответ на их Open question~2).

Дополнительно, исчерпывающим перебором всех $1\,424\,115\,231$ подрешёток
индекса $140$ решётки $A_5^*$ установлено, что явная конструкция $C_5$ из
\cite{ABPR} оптимальна на этом индексе и по ширине интервала:
$\max d=\sqrt{39/35}$.

Наконец, в $\R^3$ оптимизация метрики по семейству деформаций объёмноцентрированной
решётки улучшает таблицу работы \cite{Ivanov2011}; приведём новые значения
(каждая строка --- явная решётка и подрешётка, доступные в данных автора):
\[
\begin{array}{c|ccccccccc}
k & 16 & 20 & 21 & 22 & 23 & 24 & 26 & 28 & 30\\
\hline
d & 1{,}0254 & 1{,}1375 & 1{,}1801 & 1{,}1893 & 1{,}2140 & 1{,}3158 &
1{,}2637 & 1{,}4596 & 1{,}4730
\end{array}
\]
(для сравнения: в \cite{Ivanov2011} $k=21\mapsto1{,}133$, $k=23\mapsto1{,}137$,
$k=24\mapsto1{,}303$; значения при $k=16$, $20$, $22$, $26$, $28$--$31$ ранее
не публиковались).

\section{Открытые вопросы}

\begin{enumerate}
\item Является ли $48$ минимальным числом цветов решёточной раскраски $\R^4$?
Наши эксперименты ($\max d\approx0{,}973$ при $k=47$ и $\approx0{,}981$ при
$k=46$ по той же оптимизационной схеме) позволяют предположить утвердительный
ответ.
\item Найти аналитическое описание оптимальной решётки для $k=48$
(предельной точки максимизации $d$): числовой оптимум $1{,}04327$ подсказывает
существование выделенной алгебраической формы.
\item Останется ли значение $140$ минимальным в $\R^5$ при оптимизации по
генерическим решёткам? Аналогичный вопрос для $343$ в $\R^6$ (наш
рандомизированный поиск по методике \cite{ABPR} улучшений не дал).
\item Доказать равенство $D((3+\omega)\Lambda)^2=\tfrac73\lambda_1^2$ для
всех эйзенштейновских решёток (или описать область его справедливости).
\end{enumerate}

\subsection*{Благодарности}
Вычисления выполнены пакетами \texttt{voronoi4d} и \texttt{combigeo} (версия
1.1.0) автора. Автор признателен авторам работы \cite{ABPR} за ясную
постановку открытых вопросов.

\begin{thebibliography}{20}

\bibitem{ABPR}
A.~Arman, A.~Bondarenko, A.~Prymak, D.~Radchenko,
\emph{Upper bounds on chromatic number of $\E^n$ in low dimensions},
Electron. J. Combin. \textbf{31} (2024), no.~2, \#P2.35;
arXiv:2112.13438.

\bibitem{CoulsonPayne2007}
D.~Coulson, M.~S.~Payne,
\emph{A dense distance 1 excluding set in $\R^3$},
Austral. Math. Soc. Gaz. \textbf{34} (2007), no.~2, 97--102.

\bibitem{Coulson}
D.~Coulson,
\emph{A 15-colouring of 3-space omitting distance one},
Discrete Math. \textbf{256} (2002), 83--90.

\bibitem{deGrey}
A.~D.~N.~J. de~Grey,
\emph{The chromatic number of the plane is at least~5},
Geombinatorics \textbf{28} (2018), no.~1, 18--31.

\bibitem{DSMMV}
M.~Dutour~Sikiri\'c, D.~Madore, P.~Moustrou, F.~Vallentin,
\emph{Coloring the Voronoi tessellation of lattices},
J. London Math. Soc. (2) \textbf{104} (2021), 1135--1171.

\bibitem{ExooIsmailescu}
G.~Exoo, D.~Ismailescu,
\emph{On the chromatic number of $\R^4$},
Discrete Comput. Geom. \textbf{52} (2014), 416--423.

\bibitem{GarberMagazinov}
А.~И.~Гарбер, А.~Н.~Магазинов,
\emph{О гипотезе Вороного для четырёх- и пятимерных параллелоэдров},
УМН \textbf{77} (2022), №~1, 185--186.

\bibitem{Ivanov2006}
Л.~Л.~Иванов,
\emph{Оценка хроматического числа пространства $\R^4$},
УМН \textbf{61} (2006), №~5, 181--182.

\bibitem{Ivanov2011}
Л.~Л.~Иванов,
\emph{О хроматических числах $\R^2$ и $\R^3$ с интервалами запрещённых
расстояний},
Вестник РУДН, сер. Математика. Информатика. Физика (2011), №~1, 14--29.

\bibitem{LarmanRogers}
D.~G.~Larman, C.~A.~Rogers,
\emph{The realization of distances within sets in Euclidean space},
Mathematika \textbf{19} (1972), 1--24.

\bibitem{Nechushtan}
O.~Nechushtan,
\emph{On the space chromatic number},
Discrete Math. \textbf{256} (2002), 499--507.

\bibitem{RadoicicToth}
R.~Radoi\v{c}i\'c, G.~T\'oth,
\emph{Note on the chromatic number of the space},
Discrete and Computational Geometry, Algorithms Combin., vol.~25,
Springer, Berlin, 2003, 695--698.

\bibitem{Raigorodskii2000}
А.~М.~Райгородский,
\emph{О хроматическом числе пространства},
УМН \textbf{55} (2000), №~2, 147--148.

\bibitem{Soifer}
A.~Soifer,
\emph{The Mathematical Coloring Book},
Springer, New York, 2009.

\bibitem{Voronoi}
Г.~Ф.~Вороной,
\emph{Собрание сочинений}, т.~2,
Изд-во АН УССР, Киев, 1952.

\bibitem{VWZ}
F.~Vallentin, S.~Wei\ss bach, M.~C.~Zimmermann,
\emph{The chromatic number of 4-dimensional lattices},
arXiv:2407.03513 (2024).

\end{thebibliography}

\end{document}
